\documentclass[11pt,a4paper]{article}
\usepackage{amsmath,booktabs,geometry}
\geometry{margin=2.2cm}
\setlength{\parskip}{0.32em}
\setlength{\parindent}{0em}
\title{The clever covariate on the next ToT step}
\author{}
\date{2026}
\begin{document}
\maketitle

\section{When \(H\) exists}
One beat is one committed search step.
\(Y=1\) if that episode's search returns a solved state, and \(Y=0\) otherwise.
During the episode the query beat carries \(X\) and the action prior \(e\).
It does not carry \(Y\).
\(Y\) is copied onto every committed step of the episode only after the search returns.
The step that produced \(Y\) cannot read its own score.
The next episode can, because the window holds only beats whose \(Y\) is already known.

After that write-back the clever covariate is
\[
H=\frac{Y-e}{e(1-e)},
\]
where \(e\) is the stable score of the committed step, clipped to \((0.001,0.999)\).
The online anomaly score is \(|H|\).
The sign sets the adjustment: \(H>0.5\) is aggressive, \(H<-0.5\) is conservative, and the remainder is keep.
No UCB term, temperature, or prune threshold is added to this score.
\(H\) is not a child-count and it is not the terminal reward.
The terminal reward is \(Y\).
The expansion cost is the number of children scored at the step.

For \(Y\in\{0,1\}\) the two values are immediate.
A success gives \(H=1/e>0\).
A failure gives \(H=-1/(1-e)<0\).
If \(e<1/2\), then \(1/e>1/(1-e)\), so a success has the larger absolute value.
\(|H|\) is the size of the departure from the prior.
The sign is what says whether the departure was a success or a failure.
A detector that ranks beats by \(|H|\) alone will rank confident successes above failures whenever the prior is below one half.

\section{The judge, with the score unread}
The test search has three offers and keeps only the top one.
Each user scores three children.
The habit offer activates even-numbered users, so its rate is \(1/2\).
The growth offer activates a user unless the index is a multiple of \(3\).
The click offer activates only multiples of \(5\).
Forty users, drift at user 16, judge mix \(0.80\) toward the click score after that index.
Before the drift the judge selects habit.
After the drift it selects click.

On that unread stream, \(Y\) moves from \(0.500\) to \(0.167\).
That is eight fewer activations than \(0.500\times 24\).
The prior on the committed offer moves from \(0.50\) to \(0.20\), which is the stable score of the click offer.
\(|H|\) moves from \(2.00\) to \(1.875\).
Conditional on the outcome the split is sharp: after the drift, successes have mean \(|H|=5.00\) and failures have mean \(|H|=1.25\), because \(e=0.20\) gives \(1/e=5\) and \(1/(1-e)=1.25\).
The sign matches \(Y\) on all 24 post-drift steps.
Every success is aggressive and every failure is conservative.
The conservative share is the only channel that rises.
Children stay at 3, the window stays at 8 rows, and the prompt stays at 6 columns.
Writing the row does not change the argmax, so the business rate stays at the shortfall.

\section{The next user reads the sign}
The growth policy uses the judge through the first 12 users.
It switches the next user to the growth offer when the mean loss of the last four finished users exceeds the reference mean loss by \(0.15\).
A finished failure is exactly \(H<0\), so a run of high loss is a run of conservative scores.
The switch occurs at user 17, the second post-drift user.
From there the beam ranks by the stable score, and the growth offer is the one with stable score \(0.67\).

\begin{center}
\begin{tabular}{lrrrr}
\toprule
Policy & Pre \(Y\) & Post \(Y\) & Extra vs judge & Reading \\
\midrule
Sign not read & 0.500 & 0.167 & 0 & shortfall, \(-8\) \\
Sign read at user 17 & 0.500 & 0.625 & \(+11\) & increment, \(+3\) \\
\bottomrule
\end{tabular}
\end{center}

The second line activates 11 more of the 24 post-drift users than the drifted judge, and 3 more than the pre-drift habit rate \(0.500\times 24\).
Post-drift \(Y=0.625\) is above \(0.500\), so the business outcome itself clears its own reference.
Children remain 3 per user, so the search cost does not buy the increment.
The prior on the selected offer moves from \(0.50\) to \(0.650\).
The conservative share moves from \(0.50\) to \(0.375\), because more users activate and those steps have \(H>0\).
\(|H|\) stays at \(1.995\), against \(2.00\) before the drift.
The absolute residual moves from \(0.500\) to \(0.506\).
The anomaly magnitude is not the channel that records the increment.
The channels that move with the read are the activation rate, the prior of the chosen offer, and the share of conservative steps.

\section{What the read does not do}
On Game of 24 the same \(H\) was written into the sliding window and the next argmax still followed the frozen forest.
\(Y\) fell from \(1.00\) to \(0.25\).
The absolute residual, the conservative share, and the last-step absolute error each rose by \(0.75\), and the hop rate rose by \(0.25\).
\(|H|\) fell from \(2.38\) to \(1.63\), for the same reason as above: the prior stayed below \(1/2\), and failures have the smaller absolute score.
Scored children stayed near \(115\) per episode.
That is the control.
The formula, the window, and the anomaly score can all be present while \(Y\) does not recover.
The increment on the offer search appears when the next episode's ranking consumes the sign.
It does not appear from printing \(H\).

\end{document}
